\documentclass[a4paper]{article}
\usepackage[warn]{mathtext}
\usepackage{graphicx} % Required for inserting images
\usepackage{wrapfig}
\usepackage{enumitem}
\usepackage[center]{caption}
\usepackage[english, russian]{babel}
\usepackage{amsmath}
\usepackage{geometry}
\setcounter{totalnumber}{5}

\geometry{verbose, a4paper, tmargin=2cm, bmargin=2cm, lmargin=2.0cm, rmargin=2.0cm}

\begin{document}

\begin{titlepage}
	\centering
	{\scshape\Large МОСКОВСКИЙ ФИЗИКО-ТЕХНИЧЕСКИЙ ИНСТИТУТ \\
	(НАЦИОНАЛЬНЫЙ ИССЛЕДОВАТЕЛЬСКИЙ УНИВЕРСИТЕТ)\\ % название ВУЗа большим шрифтом
    Физтех-школа аэрокосмических технологий}
	
	\vspace{4cm} % отступ 4 см по вертикали
	{\LARGE Отчет о выполнении лабораторной работы 1.2.2}
	\vspace{1cm} % ещё 1 см
	
{\huge\bf Экспериментальная проверка \\ закона вращательного движения  \\ на крестообразном маятнике}
	
	\vspace{1cm} % ещё 1 см
	\vfill % заполнение по вертикали пустотой (LaTeX сам решит, сколько здесь отступить)
	
    \begin{flushright} % Этот текст будет с правого края страницы
	   {\LARGE Губарев Никита, Б03-502}
    \end{flushright}
    \vfill 
	\today % Дата сборки документа
\end{titlepage}
\tableofcontents 
\section{Аннотация}
Цель работы: экспериментально проверить уравнение вращательного движения тела вокруг закрепленной оси, получить зависимость углового ускорения от момента инерции и момента прикладываемых к системе сил; проанализировать влияние сил трения, действующих в оси вращения; определить момент инерции маятника.

\section{Теоретические сведения}
Уравнение вращательного движения тела вокруг закрепленной оси:
\begin{equation}
I\ddot{\varphi} = M,
\end{equation} где $\ddot{\varphi} = \dot{\omega} = \beta$ -- угловое ускорение ($\omega$ -- угловая скорость), $I$ -- полный
момент инерции тела относительно оси вращения, $M$ -- суммарный момент внешних сил относительной этой оси. \\ \\
Момент силы натяжения нити:
\begin{equation}
M_\text{н} = m_\text{н}r(g-\beta r),    
\end{equation}
где  $m_\text{н} = m_\text{п} + m_\text{г}$ -- масса платформы с перегрузком. \\ \\
Таким образом, с учетом (2) уравнение (1) может быть записано как:
\begin{equation}
(I + m_\text{н}r^2)\beta = m_\text{н}gr - M_\text{тр},
\end{equation} где $M_\text{тр}$ --  момент силы трения в оси вращения. \\ \\
В проведенных опытах $m_\text{н} r^2 << I$,  поэтому можно считать, что маятник раскручивается с постоянным угловым ускорением:
\begin{equation}
    \beta_0 = \frac{1}{I}m_\text{н}gr - \frac{M_{\text{тр}}}{I}
    \label{eq:beta_0}
\end{equation}
Т. к. грузы имеют форму полых цилиндров с внутренним и внешним радиусами и образующей $a_1$ и $a_2$ и  $h$ соответственно, момент инерции системы $I$ вычисляется при помощи теоремы Гюйгенса-Штейнера:
% Requires: \usepackage{amsmath}
\begin{equation}
    I = I_0 + \sum_{i=1}^{4} \left( \frac{1}{12} m_i h^2 + \frac{1}{4} m_i (a_1^2 + a_2^2) + m_i R_i^2 \right).
    \label{eq:moment_of_inertia}
\end{equation}

\section{Методика измерений}
\begin{enumerate}
\item Установить грузы $m_i$ на оси шкива так, чтобы маятник оказался в состоянии безразличного равновесия. Измерить положение грузов $R_i$.
\item Оценить момент силы $M_0$ трения покоя в подшипниках: намотать на меньший из шкивов нить в один слой и подвесить на нее пустую платформу; нагрузить платформу так, чтобы маятник пришел в движение. Если маятник приходит во вращение без перегрузков, занчит момент сил мал $M_0<m_пgr$, можно переходить к следующему пункту.
\item Включить компьютер и запустить программу "Kinematic", ознакомиться с инструкцией по работе с программой.
\item Намотать нить в один слой на больший из шкивов и поместить перегрузок ($m_г\approx 50 г$) на платформу. Провести опыт: с помощью программы измерьте зависимость угла поворота маятника от времени в процессе опускания платформы из верхнего в нижнее положение.

По завершении опыта перейти в раздел "Определение углового ускорения". На экране представлен график зависимости углового ускорения от угловой скорости  $\beta(\omega)$. Программа проводит через экспериментальные точки прямую методом наименьших квадратов (МНК). Записать полученные коэффициенты $\beta_0$ и $k$ прямой $\beta=\beta_0+k\omega$ и их погрешности (рассчитанные также по МНК).

Из рассмотрения следует исключить несколько начальных точек (поскольку для них существенно влияние силы трения покоя), а также, возможно, конечные точки, соответствующие изменению направления движения платформы. При наличии резких выбросов (ошибки срабатывания датчика) опыт следует переделать. 

Если на полученном графике существенны колебания зависимости $\beta(\omega)$ относительно прямой, уточните балансировку и повторите измерение. Если дальнейшее уточнение балансировки затруднено, оцените (грубо) погрешность, вносимую этими колебаниями в результат опыта.
\item Для оценки случайной погрешности опыта провести 6–8 измерений по п. 4 для фиксированных значений массы и момента инерции маятника. По результатам оценить случайную ошибку $\sigma_\beta$.
\item Провести опыт п. 4. для 6–8 значений момента силы натяжения нити, использовать перегрузки $m$ в диапазоне от 20 до 200 г на разных шкивах.
\item Измерить зависимость углового ускорения от момента инерции системы. Для этого при одном из значений массы перегрузка из предыдущего пункта ($m_г\approx 100 г$) провести измерения $\beta_0$ и $k$ при 4–5 различных значениях расстояния $R$ от оси системы до центров масс грузов. После перемещения грузов проводить балансировку.
\item Снять грузы со стержней и провести несколько опытов по опусканию грузов для определения момента инерции пустой конструкции $I_0$.
\item По результатам п. 6 проверить справедливость формулы (3). Для этого построить график зависимости начального углового ускорения $\beta_0(M_н)$ от момента силы натяжения нити (2). Убедиться, что  экспериментальные точки ложатся на прямую линию. По наклону прямой определить момент инерции $I$, а по пересечению с осью абсцисс - минимальную силу трения $M_0$. Сравнить $M_0$ с измеренным в п. 2. Оценить погрешности результатов.
\item Пользуясь формулой (3) и величиной $M_0$ , полученной в п. 9, рассчитать по результатам п. 7 моменты инерции $I$ системы при различных $R$. Оценить погрешность измерения $I$. Построить график зависимости $I(R^2)$. Использовуя график и формулу (5) определить момент инерции конструкции без грузов $I_0$ и сравнить его с измеренным в п. 8.
\end{enumerate}

\section{Результаты измерений и обработка данных}


\section{Выводы}

\newpage
\section{Приложения}
Приложение 1:
\begin{wrapfigure}{R}{0.4\textwidth}
    \includegraphics[width=1\linewidth]{/home/nik/Desktop/mipt_nik/Физос/Лабы/1.2.2/oberbek.PNG}
    \caption{Маятник Обербека}
\end{wrapfigure}
\end{document}
